\documentclass[../../../include/open-logic-section]{subfiles}

\begin{document}

\olfileid{sth}{spine}{recursion}
\olsection{قضیه‌های بازگشت ترامتناهی}

نخست باید تأیید کنیم که~\olref[valpha]{defValphas} تعریفی موفق است.
به‌طور کلی‌تر، باید اثبات کنیم هر کوششی برای ارائهٔ تعریفی با بازگشت
(ترامتناهی) به نتیجه می‌رسد. هدف این بخش همین است.

\emph{هشدار: این مطالب دشوارند}. بااین‌همه، پیام کلی بسیار ساده است:
استقرای ترامتناهی همراه با جایگزینی، مشروعیت چندین صورت از بازگشت
ترامتناهی را تضمین می‌کند.\footnote{یادآوری: همهٔ فرمول‌ها و ترم‌ها
می‌توانند پارامتر داشته باشند، مگر آنکه خلافش تصریح شود.}

\begin{defn}
	فرض کنید~$\tau(x)$ یک ترم، $f$ یک تابع و~$\alpha$ یک عدد ترتیبی باشد.
	می‌گوییم~$f$ یک \emph{تقریبِ $\alpha$ برای~$\tau$} است اگر و تنها اگر
	هم~$\dom{f}=\alpha$ و هم
	$(\forall \beta \in \alpha)f(\beta)=\tau(\funrestrictionto{f}{\beta})$.
\end{defn}

\begin{lem}[بازگشت کراندار]\ollabel{transrecursionfun}
برای هر ترم~$\tau(x)$ و هر عدد ترتیبی~$\alpha$، تقریبِ $\alpha$ یکتایی
برای~$\tau$ وجود دارد.
\end{lem}
\begin{proof}
نشان می‌دهیم برای هر~$\gamma\leq\alpha$، یک تقریبِ $\gamma$ یکتا وجود
دارد.

ابتدا یکتایی را اثبات می‌کنیم. فرض کنید~$g$ و~$h$ به‌ترتیب تقریب‌های
$\gamma$ و~$\delta$ باشند. استقرای ترامتناهی روی آرگومان‌های آن‌ها نشان
می‌دهد برای هر~$\beta\in\dom{g}\cap\dom{h}=\gamma\cap\delta=
\min(\gamma,\delta)$، داریم~$g(\beta)=h(\beta)$. پس تقریب‌های ما، اگر
وجود داشته باشند، یکتا هستند و در همهٔ مقادیر مشترک توافق دارند.

برای اثبات وجود، اکنون استقرای ترامتناهی ساده‌ای
(\olref[ordinals][opps]{simpletransrecursion}) روی اعداد ترتیبی
$\delta\leq\alpha$ به کار می‌بریم.

تابع تهی بدیهی است که یک تقریبِ $\emptyset$ است.

اگر~$g$ یک تقریبِ $\gamma$ باشد، آنگاه
$g\cup\{\tuple{\gamma,\tau(g)}\}$ یک تقریبِ
$\ordsucc{\gamma}$ است.

اگر~$\gamma$ عدد ترتیبی حدی باشد و برای هر~$\delta<\gamma$،
$g_\delta$ تقریبِ $\delta$ باشد، بگذارید
$g=\bigcup_{\delta\in\gamma}g_\delta$. این یک تابع است، زیرا
$g_\delta$های گوناگون ما در همهٔ مقادیر مشترک توافق دارند. همچنین اگر
$\delta\in\gamma$، آنگاه
$g(\delta)=g_{\ordsucc{\delta}}(\delta)=
\tau(\funrestrictionto{g_{\ordsucc{\delta}}}{\delta})=
\tau(\funrestrictionto{g}{\delta})$.

این، اثبات با استقرای ترامتناهی را کامل می‌کند.
\end{proof}

اگر به خود اجازه دهیم به‌جای یک تابع، یک \emph{ترم} تعریف کنیم، می‌توانیم
کران~$\alpha$ را از نتیجهٔ پیشین حذف کنیم. در بیان و اثبات نتیجهٔ زیر،
هرگاه~$\sigma$ یک ترم باشد، می‌گذاریم
$\funrestrictionto{\sigma}{\alpha}=
\Setabs{\tuple{\beta,\sigma(\beta)}}{\beta\in\alpha}$.

\begin{thm}[بازگشت عمومی]\ollabel{transrecursionschema}
برای هر ترم~$\tau(x)$ می‌توانیم ترمی مانند~$\sigma(x)$ را به‌صراحت تعریف
کنیم، به‌طوری‌که برای هر عدد ترتیبی~$\alpha$،
$\sigma(\alpha)=\tau(\funrestrictionto{\sigma}{\alpha})$.
\end{thm}

\begin{proof}
برای هر~$\alpha$، بنا بر~\olref{transrecursionfun} تقریبِ $\alpha$ یکتایی
مانند~$f_\alpha$ برای~$\tau$ وجود دارد. $\sigma(\alpha)$ را
$f_{\ordsucc{\alpha}}(\alpha)$ تعریف کنید. اکنون:
	\begin{align*}
		\sigma(\alpha) &=
		f_{\ordsucc{\alpha}}(\alpha) \\&=
		\tau(\funrestrictionto{f_{\ordsucc{\alpha}}}{\alpha}) \\&=
		\tau(\Setabs{\tuple{\beta, f_{\ordsucc{\alpha}}(\beta)}}{\beta \in \alpha}) \\&=
		\tau(\Setabs{\tuple{\beta, f_{\ordsucc{\beta}}(\beta)}}{\beta \in \alpha})\\&=
		\tau(\funrestrictionto{\sigma}{\alpha})
	\end{align*}
زیرا، همان‌گونه که در~\olref{transrecursionfun} دیدیم، برای همهٔ
$\beta<\alpha$ داریم
$f_{\ordsucc{\beta}}(\beta)=f_{\ordsucc{\alpha}}(\beta)$.
\end{proof}
\noindent
توجه کنید که~\olref{transrecursionschema} یک \emph{طرحواره} است. نکتهٔ
اساسی این است که نمی‌توانیم انتظار داشته باشیم~$\sigma$ یک تابع، یعنی
نوع معینی از \emph{مجموعه}، را تعریف کند؛ زیرا در آن صورت~$\dom{\sigma}$
مجموعهٔ همهٔ اعداد ترتیبی می‌بود و این با پارادوکس بورالی--فورتی
(\olref[ordinals][basic]{buraliforti}) تناقض داشت.

بااین‌حال، هنوز باید نشان دهیم~\olref{transrecursionschema} تعریف ما از
$V_\alpha$ها را موجه می‌سازد. شاید این امر بی‌درنگ آشکار نباشد، اما با
آخرین صورت سادهٔ بازگشت ترامتناهی روشن خواهد شد.

\begin{thm}[بازگشت ساده]\ollabel{simplerecursionschema}
برای هر دو ترم~$\tau(x)$ و~$\theta(x)$ و هر مجموعهٔ~$A$، می‌توانیم ترمی
مانند~$\sigma(x)$ را به‌صراحت تعریف کنیم به‌طوری‌که:
\begin{align*}
	\sigma(\emptyset) &= A\\
	\sigma(\ordsucc{\alpha}) &= \tau(\sigma(\alpha)) &&
		\text{برای هر عدد ترتیبی }\alpha\\
	\sigma(\alpha) &= \theta(\ran{\funrestrictionto{\sigma}{\alpha}})&&
	\text{اگر }\alpha\text{ عدد ترتیبی حدی باشد}
\end{align*}
\end{thm}

\begin{proof}
ابتدا ترم~$\xi(x)$ را چنین تعریف می‌کنیم:
%t اگر $x=\emptyset$ باشد یا $x$ تابعی نباشد که دامنه‌اش یک عدد ترتیبی است، $\xi(x)=A$؛
%در غیر این صورت بگذارید:
\[
	\xi(x) =
	\begin{cases}
			A &\text{اگر $x=\emptyset$ باشد یا $x$ تابعی نباشد که}\\
			&\hspace{1em}\text{دامنه‌اش یک عدد ترتیبی است؛ در غیر این صورت:}\\
			\tau(x(\alpha)) & \text{اگر $\dom{x}=\ordsucc{\alpha}$}\\
			\theta(\ran{x}) & \text{اگر $\dom{x}$ یک عدد ترتیبی حدی باشد}
	\end{cases}
\]
بنا بر~\olref{transrecursionschema}، ترمی مانند~$\sigma(x)$ وجود دارد
به‌طوری‌که برای هر عدد ترتیبی~$\alpha$،
$\sigma(\alpha)=\xi(\funrestrictionto{\sigma}{\alpha})$؛ افزون بر این،
$\funrestrictionto{\sigma}{\alpha}$ تابعی با دامنهٔ~$\alpha$ است. با
استقرای ترامتناهی ساده
(\olref[ordinals][opps]{simpletransrecursion}) نشان می‌دهیم~$\sigma$
ویژگی‌های لازم را دارد.

نخست، $\sigma(\emptyset)=\xi(\emptyset)=A$.

سپس، $\sigma(\ordsucc{\alpha})=
\xi(\funrestrictionto{\sigma}{\ordsucc{\alpha}})=
\tau(\funrestrictionto{\sigma}{\ordsucc{\alpha}}(\alpha))=
\tau(\sigma(\alpha))$.

سرانجام، اگر~$\alpha$ حدی باشد،
$\sigma(\alpha)=\xi(\funrestrictionto{\sigma}{\alpha})=
\theta(\ran{\funrestrictionto{\sigma}{\alpha}})$.
\end{proof}
\noindent
اکنون برای موجه‌ساختن~\olref[valpha]{defValphas} کافی است
$A=\emptyset$، $\tau(x)=\Pow{x}$ و~$\theta(x)=\bigcup x$ را اختیار
کنیم. سرانجام تعریف~$V_\alpha$ها موجه شد!{}


\end{document}
